\ifdefined\gkver\else\def\gkver{student}\fi
\documentclass[\gkver]{gaokaozhenti}
\begin{document}

\chapter{2014年北京理综}
%% paperType: 理综物理部分

\begin{enumerate}
\item
%% number: 13
%% paperName: 2014年北京理综第 13 题 6 分
%% typeId: 1
%% type: 单选题
%% chapter: 气体性质
%% point: 物体的内能
%% method: 无
%% score: 6
%% degree: 650
%% duplicateId: 0
%% body:
下列说法中正确的是\xzanswer{B}
\fourchoices[answer=B]
{物体温度降低，其分子热运动的平均动能增大}
{物体温度升高，其分子热运动的平均动能增大}
{物体温度降低，其内能一定增大}
{物体温度不变，其内能一定不变}
%% answer: B
%% memo:
\memoanswer{
温度是分子平均动能的标志，温度降低分子热运动的平均动能减小，反之增大，A 项错误；B 项正确；物体的内能包括所有分子的动能和势能之和，温度降低，分子动能减小但是分子势能不能确定，所以内能不能确定，CD 项错误。
}

\item
%% number: 14
%% paperName: 2014年北京理综第 14 题 6 分
%% typeId: 1
%% type: 单选题
%% chapter: 原子和原子核
%% point: 质能方程
%% method: 无
%% score: 6
%% degree: 650
%% duplicateId: 0
%% body:
质子、中子和氘核的质量分别为 $m_{1}$、$m_{2}$ 和 $m_{3}$。当一个质子和一个中子结合成氘核时，释放的能量是（$c$ 表示真空中的光速）\xzanswer{C}
\fourchoices[answer=C]
{$(m_{1}+m_{2}-m_{3})c$}
{$(m_{1}-m_{2}-m_{3})c$}
{$(m_{1}+m_{2}-m_{3})c^{2}$}
{$(m_{1}-m_{2}-m_{3})c^{2}$}
%% answer: C
%% memo:
\memoanswer{
原子核反应经过裂变或聚变时，质量会发生亏损，损失的能量以能量的形式释放，根据质能方程可知释放的能量为：$\Delta E=(m_{1}+m_{2}-m_{3})c^{2}$，C 项正确。
}

\item
%% number: 15
%% paperName: 2014年北京理综第 15 题 6 分
%% typeId: 1
%% type: 单选题
%% chapter: 电场
%% point: 电场线
%% method: 无
%% score: 6
%% degree: 650
%% duplicateId: 0
%% body:
如图所示，实线表示某静电场的电场线，虚线表示该电场的等势面。下列判断正确的是\xzanswer{D}
\onepicture[w=3.4cm]{figs/15.png}
\fourchoices[answer=D]
{1、2 两点的场强相等}
{1、3 两点的场强相等}
{1、2 两点的电势相等}
{2、3 两点的电势相等}
%% answer: D
%% memo:
\memoanswer{
电场线的疏密程度反映了电场强度的大小，所以 1 点的场强大于 2 点；电场线的切线方向表示场强的方向，2 点和 3 点的场强方向不同，AB 项错误；沿电场线电势降低，1 点电势高于 2 点电势，CD 项错误；2 点和 3 点电势相等，D 项正确。
}

\item
%% number: 16
%% paperName: 2014年北京理综第 16 题 6 分
%% typeId: 1
%% type: 单选题
%% chapter: 磁场
%% point: 带电粒子在磁场中的运动
%% method: 无
%% score: 6
%% degree: 450
%% duplicateId: 0
%% body:
带电粒子 $a$、$b$ 在同一匀强磁场中做匀速圆周运动，它们的动量大小相等，$a$ 运动的半径大于 $b$ 运动的半径。若 $a$、$b$ 的电荷量分别为 $q_{a}$、$q_{b}$，质量分别为 $m_{a}$、$m_{b}$，周期分别为 $T_{a}$、$T_{b}$。则一定有\xzanswer{A}
\fourchoices[answer=A]
{$q_{a}<q_{b}$}
{$m_{a}<m_{b}$}
{$T_{a}<T_{b}$}
{$\frac{q_{a}}{m_{a}}<\frac{q_{b}}{m_{b}}$}
%% answer: A
%% memo:
\memoanswer{
带电粒子在匀强磁场中做匀速圆周运动，由牛顿第二定律有：$qvB=m\frac{v^{2}}{r}$ 得 $r=\frac{mv}{qB}$，因为两个粒子的动量相等，且 $r_{a}>r_{b}$，所以 $q_{a}<q_{b}$，A 项正确；速度不知道，所以质量关系不确定，B 项错误；又因为 $T=\frac{2\pi m}{qB}$，质量关系不知道，所以周期关系不确定，CD 项错误。
}

\item
%% number: 17
%% paperName: 2014年北京理综第 17 题 6 分
%% typeId: 1
%% type: 单选题
%% chapter: 机械振动
%% point: 振动图象
%% method: 无
%% score: 6
%% degree: 450
%% duplicateId: 0
%% body:
一简谐机械横波沿 $x$ 轴正方向传播，波长为 $\lambda$，周期为 $T$。$t=0$ 时刻的波形如图~\ref{2014北京理综17a} 所示，$a$、$b$ 是波上的两个质点。图~\ref{2014北京理综17b} 是波上某一质点的振动图像。下列说法中正确的是\xzanswer{D}
\twopicture[labelA=2014北京理综17a, labelB=2014北京理综17b, numA=1, numB=2, numstyle=arabic, wA=5.3cm, wB=5.3cm]
{figs/17a.png}
{figs/17b.png}
\fourchoices[answer=D]
{$t=0$ 时质点 $a$ 的速度比质点 $b$ 的大}
{$t=0$ 时质点 $a$ 的加速度比质点 $b$ 的小}
{图~\ref{2014北京理综17b} 可以表示质点 $a$ 的振动}
{图~\ref{2014北京理综17b} 可以表示质点 $b$ 的振动}
%% answer: D
%% memo:
\memoanswer{
由图 1 的波形图可知 $t=0$ 时刻，$a$ 在波峰，速度为零，加速度最大；$b$ 在平衡位置，速度最大，加速度为零，AB 项错误；根据微平移法可以确定 0 时刻 $b$ 质点向下振动，所以图 2 是质点 $b$ 的振动图象，C 项错误，D 项正确。
}

\item
%% number: 18
%% paperName: 2014年北京理综第 18 题 6 分
%% typeId: 1
%% type: 单选题
%% chapter: 运动和力
%% point: 超重失重
%% method: 无
%% score: 6
%% degree: 450
%% duplicateId: 0
%% body:
应用物理知识分析生活中的常见现象，可以使物理学习更加有趣和深入。例如平伸手掌托起物体，由静止开始竖直向上运动，直至将物体抛出。对此现象分析正确的是\xzanswer{D}
\fourchoices[answer=D]
{手托物体向上运动的过程中，物体始终处于超重状态}
{手托物体向上运动的过程中，物体始终处于失重状态}
{在物体离开手的瞬间，物体的加速度大于重力加速度}
{在物体离开手的瞬间，手的加速度大于重力加速度}
%% answer: D
%% memo:
\memoanswer{
物体具有向上的加速度为超重，向下的加速度为失重，手托物体抛出的过程，必定有一段加速过程，即超重过程，从加速后到手和物体分离的过程中，可以匀速也可以减速，因此可能失重，也可能既不超重也不失重，A、B 错误。手与物体分离时的力学条件为：手与物体之间的压力 $N=0$，分离后手和物体一定减速，物体减速的加速度为 $g$，手减速要比物体快才会分离，因此手的加速度大于 $g$，C 错误，D 正确。
}

\item
%% number: 19
%% paperName: 2014年北京理综第 19 题 6 分
%% typeId: 1
%% type: 单选题
%% chapter: 运动和力
%% point: 牛顿第一定律
%% method: 无
%% score: 6
%% degree: 650
%% duplicateId: 0
%% body:
伽利略创造的把实验、假设和逻辑推理相结合的科学方法，有力地促进了人类科学认识的发展。利用如图所示的装置做如下实验：小球从左侧斜面上的 $O$ 点由静止释放后沿斜面向下运动，并沿右侧斜面上升。斜面上先后铺垫三种粗糙程度逐渐减低的材料时，小球沿右侧斜面上升到的最高位置依次为 1、2、3。根据三次实验结果的对比，可以得到的最直接的结论是\xzanswer{A}
\onepicture[w=7.6cm]{figs/19.png}
\fourchoices[answer=A]
{如果斜面光滑，小球将上升到与 $O$ 点等高的位置}
{如果小球不受力，它将一直保持匀速运动或静止状态}
{如果小球受到力的作用，它的运动状态将发生改变}
{小球受到的力一定时，质量越大，它的加速度越小}
%% answer: A
%% memo:
\memoanswer{
由题给出斜面上铺垫三种粗糙程度逐渐降低的材料，小球的位置逐渐升高，不难想象，当斜面绝对光滑时，小球在斜面上运动没有能量损失，可以上升到与 $O$ 点等高的位置，这是可以得到的直接结论，A 正确，B、C、D 尽管也正确，但不是本实验得到的直接结论，故错误。
}

\item
%% number: 20
%% paperName: 2014年北京理综第 20 题 6 分
%% typeId: 1
%% type: 单选题
%% chapter: 光学
%% point: 光的折射
%% method: 无
%% score: 6
%% degree: 450
%% duplicateId: 0
%% body:
以往，已知材料的折射率都为正值（$n>0$）。现已有针对某些电磁波设计制作的人工材料，其折射率可以为负值（$n<0$），称为负折射率材料。位于空气中的这类材料，入射角 $i$ 与折射角 $r$ 依然满足 $\frac{\sin i}{\sin r}=n$，但是折射线与入射线位于法线的同一侧（此时折射角取负值）。现空气中有一上下表面平行的负折射率材料，一束电磁波从其上表面射入，下表面射出。若该材料对此电磁波的折射率 $n=-1$，正确反映电磁波穿过该材料的传播路径的示意图是\xzanswer{B}
\fourchoices[answer=B, ispicture=true, w=3.3cm]
{figs/20a.png}
{figs/20b.png}
{figs/20c.png}
{figs/20d.png}
%% answer: B
%% memo:
\memoanswer{
光从空气射入到材料中时，根据题意可知入射光线和折射光线位于法线的同侧，A 项错误；根据折射定律可知：$\frac{\sin i}{\sin r}=-1$，可知折射角等于入射角，C 项错误；根据光路可逆，光从材料射入到空气时，同理入射光线和折射光线也位于法线的同侧，D 项错误，入射角也等于折射角，B 项正确。
}

\item
%% number: 21
%% paperName: 2014年北京理综第 21 题 18 分
%% typeId: 4
%% type: 实验题
%% chapter: 电路
%% point: 电路实验
%% method: 无
%% score: 18
%% degree: 450
%% duplicateId: 0
%% body:
利用电流表和电压表测定一节干电池的电动势和内电阻。要求尽量减小实验误差。
\begin{enumerate}
	\item
	应该选择的实验电路是\tkanswer{甲}（选填“甲”或“乙”）。

	\twopicture[labelA=2014北京理综21a, labelB=2014北京理综21b, numA=甲, numB=乙, labelprefix={}, wA=6.0cm, wB=6.0cm]
	{figs/21a.png}
	{figs/21b.png}
	\item
	现有电流表（$0\sim0.6\UA$）、开关和导线若干，以及以下器材：

	A．电压表（$0\sim15\UV$）

	B．电压表（$0\sim3\UV$）

	C．滑动变阻器（$0\sim50\UO$）

	D．滑动变阻器（$0\sim500\UO$）

	实验中电压表应选用\tkanswer{B}；滑动变阻器应选用\tkanswer{C}。（选填相应器材前的字母）
	\item
	某位同学记录的 6 组数据如下表所示，其中 5 组数据的对应点已经标在图~\ref{2014北京理综21c} 的坐标纸上，请标出余下一组数据的对应点，并画出 $U-I$ 图线。

	\begin{center}
	\begin{tabular}{|c|c|c|c|c|c|c|}
	\hline
	序号 & 1 & 2 & 3 & 4 & 5 & 6 \\
	\hline
	电压 $U$（V） & 1.45 & 1.40 & 1.30 & 1.25 & 1.20 & 1.10 \\
	\hline
	电流 $I$（A） & 0.060 & 0.120 & 0.240 & 0.260 & 0.360 & 0.480 \\
	\hline
	\end{tabular}
	\end{center}

	\onepicture[num=2, label=2014北京理综21c, numstyle=arabic, w=8.0cm]{figs/21c.png}
	\item
	根据（3）中所画图线可得出干电池的电动势 $E=$\tkanswer{1.50} $\UV$，内电阻 $r=$\tkanswer{0.83} $\UO$。
	\item
	实验中，随着滑动变阻器滑片的移动，电压表的示数 $U$ 及干电池的输出功率 $P$ 都会发生变化。图 3 的各示意图中正确反映 $P-U$ 关系的是\tkanswer{C}。
\end{enumerate}

\fourpicture[labelA=2014北京理综21d, labelB=2014北京理综21e, labelC=2014北京理综21f, labelD=2014北京理综21g, numA=A, numB=B, numC=C, numD=D, labelprefix={}, w=2.9cm]
{figs/21d.png}
{figs/21e.png}
{figs/21f.png}
{figs/21g.png}

%% answer:
\jdanswer{
\begin{enumerate}
	\item
	甲
	\item
	BC
	\item
	$U-I$ 图线见答案图。
	\onepicture[w=7.5cm]{figs/21_an.png}
	\item
	1.50（1.49$\sim$1.51），0.83（0.81$\sim$0.85）
	\item
	C
\end{enumerate}
}
%% memo:
\memoanswer{
（1）根据 $U=E-Ir$ 测量电源电动势和内阻时，需要测出多组对应的路端电压 $U$ 和干路电流 $I$，电压表和电流表内阻影响会造成实验误差。电源内阻较小，所以电流表分压影响较大，因此应选择甲电路。

（3）$U-I$ 图线见答案图。

（4）根据 $U-I$ 图像，电源的电动势等于纵轴的截距，内阻为斜率的绝对值。

（5）电源输出功率 $P=UI=U\left(\frac{E-U}{r}\right)=-\frac{1}{r}U^{2}+\frac{E}{r}U$，$P-U$ 图像为开口向下的二次函数，应选（C）。
}

\item
%% number: 22
%% paperName: 2014年北京理综第 22 题 16 分
%% typeId: 6
%% type: 计算题
%% chapter: 运动和力
%% point: 动量守恒定律
%% method: 无
%% score: 16
%% degree: 450
%% duplicateId: 0
%% body:
如图所示，竖直平面内的四分之一圆弧轨道下端与水平桌面相切，小滑块 A 和 B 分别静止在圆弧轨道的最高点和最低点。现将 A 无初速度释放，A 与 B 碰撞后结合为一个整体，并沿桌面滑动。已知圆弧轨道光滑，半径 $R=0.2\Um$；A 和 B 的质量相等；A 和 B 整体与桌面之间的动摩擦因数 $\mu=0.2$。取重力加速度 $g=10\Umsq$。求：
\begin{enumerate}
	\item
	碰撞前瞬间 A 的速率 $v$；
	\item
	碰撞后瞬间 A 和 B 整体的速率 $v'$；
	\item
	A 和 B 整体在桌面上滑动的距离 $l$。
\end{enumerate}
\onepicture[align=right, w=6.5cm]{figs/22.png}
%% answer:
\jdanswer{
\begin{enumerate}
	\item
	$2\Ums$
	\item
	$1\Ums$
	\item
	$0.25\Um$
\end{enumerate}
}
%% memo:
\memoanswer{
（1）从圆弧最高点滑到最低点的过程中机械能守恒则有：$\frac{1}{2}m_{A}v_{A}^{2}=m_{A}gR$，可以得出 $v_{A}=\sqrt{2gR}=2\Ums$。

（2）在圆弧底部和滑块 B 相撞，动量守恒：$(m_{A}+m_{B})v'=m_{A}v_{A}$，得 $v'=\frac{1}{2}v_{A}=1\Ums$。

（3）滑块 A 和 B 粘在一起在桌面上滑行，摩擦力做负功，由动能定理可得：$f\cdot l=\frac{1}{2}(m_{A}+m_{B})v'^{2}=\frac{1}{2}m_{A}v_{A}^{2}=\frac{1}{2}m_{A}gh$；摩擦力 $f=\mu(m_{A}+m_{B})g=\mu 2m_{A}g$；于是可以得到 $2\mu m_{A}g\cdot l=\frac{1}{2}m_{A}gR$，则 $l=\frac{R}{4\mu}=0.25\Um$。
}

\item
%% number: 23
%% paperName: 2014年北京理综第 23 题 18 分
%% typeId: 6
%% type: 计算题
%% chapter: 曲线运动
%% point: 天体运动
%% method: 无
%% score: 18
%% degree: 450
%% duplicateId: 0
%% body:
万有引力定律揭示了天体运行规律与地上物体运动规律具有内在的一致性。
\begin{enumerate}
	\item
	用弹簧秤称量一个相对于地球静止的小物体的重量，随称量位置的变化可能会有不同的结果。已知地球质量为 $M$，自转周期为 $T$，万有引力常量为 $G$。将地球视为半径为 $R$、质量均匀分布的球体，不考虑空气的影响。设在地球北极地面称量时，弹簧秤的读数是 $F_{0}$。

	a．若在北极上空高出地面 $h$ 处称量，弹簧秤读数为 $F_{1}$，求比值 $F_{1}/F_{0}$ 的表达式，并就 $h=1.0\%R$ 的情形算出具体数值（计算结果保留两位有效数字）；

	b．若在赤道地面称量，弹簧秤读数为 $F_{2}$，求比值 $F_{2}/F_{0}$ 的表达式。
	\item
	设想地球绕太阳公转的圆周轨道半径 $r$、太阳的半径 $R_{s}$ 和地球的半径 $R$ 三者均减小为现在的 1.0\%，而太阳和地球的密度均匀且不变。仅考虑太阳和地球之间的相互作用，以现实地球的 1 年为标准，计算“设想地球”的 1 年将变为多长？
\end{enumerate}
%% answer:
\jdanswer{
\begin{enumerate}
	\item
	a．$\frac{F_{1}}{F_{0}}=\frac{R^{2}}{(R+h)^{2}}$，0.98

	b．$\frac{F_{2}}{F_{0}}=1-\frac{4\pi^{2}R^{3}}{GMT^{2}}$
	\item
	1 年
\end{enumerate}
}
%% memo:
\memoanswer{
（1）设小物体质量为 $m$。

（A）在北极地面 $G\frac{Mm}{R^{2}}=F_{0}$；在北极上空高出地面 $h$ 处 $G\frac{Mm}{(R+h)^{2}}=F_{1}$，得 $\frac{F_{1}}{F_{0}}=\frac{R^{2}}{(R+h)^{2}}$；当 $h=1.0\%R$ 时 $\frac{F_{1}}{F_{0}}=\frac{1}{1.01^{2}}\approx0.98$。

（B）在赤道地面，小物体随地球自转做匀速圆周运动，受到万有引力和弹簧秤的作用力，有 $G\frac{Mm}{R^{2}}-F_{2}=m\frac{4\pi^{2}}{T^{2}}R$，得 $\frac{F_{2}}{F_{0}}=1-\frac{4\pi^{2}R^{3}}{GMT^{2}}$。

（2）地球绕太阳做匀速圆周运动，受到太阳的万有引力，设太阳质量为 $M_{s}$，地球质量为 $M$，地球公转周期为 $T_{E}$，有 $G\frac{M_{s}M}{r^{2}}=Mr\frac{4\pi^{2}}{T_{E}^{2}}$，得 $T_{E}=\sqrt{\frac{4\pi^{2}r^{3}}{GM_{s}}}=\sqrt{\frac{3\pi r^{3}}{G\rho R_{s}^{3}}}$，其中 $\rho$ 为太阳的密度。由上式可知，地球公转周期 $T_{E}$ 仅与太阳的密度、地球公转轨道半径与太阳半径之比有关。因此“设想地球”的 1 年与现实地球的 1 年时间相同。
}

\item
%% number: 24
%% paperName: 2014年北京理综第 24 题 20 分
%% typeId: 6
%% type: 计算题
%% chapter: 电磁感应
%% point: 电磁感应能量分析
%% method: 无
%% score: 20
%% degree: 250
%% duplicateId: 0
%% body:
导体切割磁感线的运动可以从宏观和微观两个角度来认识。如图所示，固定于水平面的 U 形导线框处于竖直向下的匀强磁场中，金属直导线 $MN$ 在与其垂直的水平恒力 $F$ 作用下，在导线框上以速度 $v$ 做匀速运动，速度 $v$ 与恒力 $F$ 方向相同；导线 $MN$ 始终与导线框形成闭合电路。已知导线 $MN$ 电阻为 $R$，其长度 $L$ 恰好等于平行轨道间距，磁场的磁感应强度为 $B$，忽略摩擦阻力和导线框的电阻。
\begin{enumerate}
	\item
	通过公式推导验证：在 $\Delta t$ 时间内，$F$ 对导线 $MN$ 所做的功 $W$ 等于电路获得的电能 $W_{\text{电}}$，也等于导线 $MN$ 中产生的热量 $Q$；
	\item
	若导线 $MN$ 的质量 $m=8.0\Ug$、长度 $L=0.10\Um$，感应电流 $I=1.0\UA$，假设一个原子贡献一个自由电子，计算导线 $MN$ 中电子沿导线长度方向定向移动的平均速率 $v_{e}$（下表中列出一些你可能会用到的数据）；

	\begin{center}
	\begin{tabular}{|c|c|}
	\hline
	阿伏伽德罗常数 $N_{A}$ & $6.0\times10^{23}\Umoln$ \\
	\hline
	元电荷 $e$ & $1.6\times10^{-19}\UC$ \\
	\hline
	导线 $MN$ 的摩尔质量 $\mu$ & $6.0\times10^{-2}\Ukgmol$ \\
	\hline
	\end{tabular}
	\end{center}
	\item
	经典物理学认为，金属的电阻源于定向运动的自由电子和金属离子（即金属原子失去电子后的剩余部分）的碰撞。展开你想象的翅膀，给出一个合理的自由电子的运动模型；在此基础上，求出导线 $MN$ 中金属离子对一个自由电子沿导线长度方向的平均作用力 $\overline{f}$ 的表达式。
\end{enumerate}
\onepicture[align=right, w=6.0cm]{figs/24.png}
%% answer:
\jdanswer{
\begin{enumerate}
	\item
	导线产生的感应电动势 $E=BLv$；导线匀速运动，受力平衡，$F=F_{\text{安}}=BIL$；在 $\Delta t$ 时间内，外力 $F$ 对导线做功 $W=Fv\Delta t=F_{\text{安}}v\Delta t=BILv\Delta t$；电路获得的电能 $W_{\text{电}}=qE=IE\Delta t=BILv\Delta t$；可见，$F$ 对导线 $MN$ 做的功等于电路获得的电能 $W_{\text{电}}$；导线 $MN$ 中产生的热量 $Q=I^{2}R\Delta t=I\Delta t\cdot IR=qE=W_{\text{电}}$；可见，电路获得的电能 $W_{\text{电}}$ 等于导线 $MN$ 中产生的热量 $Q$。
	\item
	$v_{e}=7.8\times10^{-6}\Ums$
	\item
	$\overline{f}=evB$
\end{enumerate}
}
%% memo:
\memoanswer{
（1）导线产生的感应电动势 $E=BLv$；导线匀速运动，受力平衡，$F=F_{\text{安}}=BIL$；在 $\Delta t$ 时间内，外力 $F$ 对导线做功 $W=Fv\Delta t=F_{\text{安}}v\Delta t=BILv\Delta t$；电路获得的电能 $W_{\text{电}}=qE=IE\Delta t=BILv\Delta t$；可见，$F$ 对导线 $MN$ 做的功等于电路获得的电能 $W_{\text{电}}$；导线 $MN$ 中产生的热量 $Q=I^{2}R\Delta t=I\Delta t\cdot IR=qE=W_{\text{电}}$；可见，电路获得的电能 $W_{\text{电}}$ 等于导线 $MN$ 中产生的热量 $Q$。

（2）导线 $MN$ 中具有的原子数为 $N=\frac{m}{\mu}N_{\mathrm{A}}$，因为一个金属原子贡献一个电子，所以导线 $MN$ 中的自由电子数也是 $N$；导线 $MN$ 单位体积内的自由电子数 $n=\frac{N}{SL}$，其中 $S$ 为导线 $MN$ 的横截面积；因为电流 $I=nv_{e}Se$，所以 $v_{\mathrm{e}}=\frac{I}{nSe}=\frac{IL}{Ne}=\frac{IL\mu}{mN_{\mathrm{A}}e}$，解得 $v_{e}=7.8\times10^{-6}\Ums$。

（3）下列解法的共同假设：所有自由电子（简称电子，下同）以同一方式运动。

方法一：动量解法。设电子在第一次碰撞结束至下一次碰撞结束之间的运动都相同，经历的时间为 $\Delta t$，电子的动量变化为零。因为导线 $MN$ 的运动，电子受到沿导线方向的洛伦兹力 $f_{\text{洛}}$ 的作用 $f_{\text{洛}}=evB$；沿导线方向，电子只受到金属离子的作用力和 $f_{\text{洛}}$ 作用，所以 $I_{f}-f_{\text{洛}}\Delta t=0$，其中 $I_{f}$ 为金属离子对电子的作用力的冲量，其平均作用力为 $f$，则 $I_{f}=f\Delta t$，得 $f=f_{\text{洛}}=evB$。

方法二：能量解法。设电子从导线的一端到达另一端经历的时间为 $t$，在这段时间内，通过导线一端的电子总数 $N=\frac{It}{e}$；电阻上产生的焦耳热是由于克服金属离子对电子的平均作用力 $f$ 做功产生的。在时间 $t$ 内，总的焦耳热 $Q=NfL$；根据能量守恒定律，有 $Q=W_{\text{电}}=EIt=BLvIt$，所以 $f=evB$。

方法三：力的平衡解法。因为电流不变，所以假设电子以速度 $v_{e}$ 相对导线做匀速直线运动。因为导线 $MN$ 的运动，电子受到沿导线方向的洛伦兹力 $f_{\text{洛}}$ 的作用 $f_{\text{洛}}=evB$；沿导线方向，电子只受到金属离子的平均作用力 $f$ 和 $f_{\text{洛}}$ 作用，二力平衡，即 $f=f_{\text{洛}}=evB$。
}

\end{enumerate}

\end{document}
